\documentclass[11pt,a4paper]{article}
\usepackage[margin=1in]{geometry}
\usepackage{amsmath,amssymb}
\usepackage[utf8]{inputenc}
% \usepackage[T1]{fontenc}  % requires full texlive
\usepackage{hyperref}

\title{A Lean-Verified Crosswalk Between Stochastic Semantics,\\Contracts, and Hoare Logic}
\author{June Kim\\{\small Independent $\cdot$ \texttt{june@june.kim}}}
\date{}

\begin{document}
\maketitle
\thispagestyle{empty}

\noindent
We present a Lean~4 artifact (about 2,800 lines, zero \texttt{sorry}, standard library only) that formalizes a stochastic information-processing pipeline and uses it to assemble a precise crosswalk across several communities in applied category theory. The core development models six stages of processing as morphisms in the Kleisli category of a lawful monad, with stage contracts represented as predicates preserved by composition. Within Lean, we prove a contract-composition theorem, derive two contracts as optimality conditions, reconstruct pigeonhole and stochasticity arguments from first principles, and formalize a recursive consolidation tower with budget-based termination.

A compiled Hoare-logic layer develops Hoare triples, weakest preconditions, and a bridge theorem showing that the original contract-preservation predicate is equivalent to a Hoare triple with trivial precondition. The same pipeline properties are proved from two independent angles---one functional, one imperative---and their equivalence is machine-checked. The crosswalk between stochastic semantics, contracts, and program logic is compiled, not just described.

\medskip\noindent\textbf{Exact identifications.}
The ambient category of the development is $\mathrm{Kleisli}(\mathcal{D}) = \mathrm{FinStoch}$, the paradigmatic Markov category~\cite{Fritz2020}. The Lean \texttt{Support} typeclass aligns with the possibilistic shadow of Fritz, Perrone, and Rezagholi~\cite{FPR2021}: \texttt{support\_bind} is the monad multiplication for $\mathrm{supp} : V \Rightarrow H$. Bonchi, Di~Lavore, Rom\'an, and Staton~\cite{BDRS2025} prove that $\mathrm{Stoch}$ is a posetal imperative category (Corollary~76) and derive standard Hoare logic rules as theorems (Theorem~79). The Lean bridge theorem identifies behavioral contracts on stochastic morphisms with classical Hoare reasoning. The Hoare layer includes a nontriviality counterexample, ruling out vacuous readings.

\medskip\noindent\textbf{Plausible correspondences.}
The contract predicate is close to a Hoare postcondition in graded Hoare logic~\cite{GKOS2021}; the handshake condition resembles the restricted pointwise order in graded-monadic program logics~\cite{KGSU2026}; and neighboring ideas appear in categorical cybernetics, including selection relations on parametrised optics~\cite{CGHR2021} and compositional equilibrium predicates~\cite{GKLF2020}.

\medskip\noindent\textbf{Open problems.}
(1)~The free $\mathbb{N}$-semimodule monad~\cite{KW2021} composed with the support monad morphism would track support cardinality through stochastic composition, connecting Leinster's magnitude~\cite{Leinster2021} to Fritz's possibilistic shadow. This composite has not been constructed.
(2)~Determinantal point processes have no categorical semantics.
(3)~Fritz's categorical entropy and data processing inequality have no accompanying achievability or rate-distortion results inside Markov categories.

\medskip\noindent\textbf{Keywords:} Markov categories, Hoare logic, Lean~4, possibilistic shadow, graded monads, imperative categories, weakest precondition, compositional contracts.

\medskip\noindent
The Lean development, vocabulary crosswalk, and full reference table are publicly available at \texttt{github.com/kimjune01/natural-framework} and \texttt{june.kim/framework-lexicon}.

\begin{thebibliography}{99}
\small

\bibitem{BDRS2025}
F.~Bonchi, E.~Di~Lavore, M.~Rom\'an, S.~Staton.
Program Logics via Distributive Monoidal Categories.
\textit{arXiv:2507.18238}, 2025.

\bibitem{CGHR2021}
M.~Capucci, B.~Gavranovi\'c, J.~Hedges, E.~Rischel.
Towards Foundations of Categorical Cybernetics.
\textit{ACT}, 2021.

\bibitem{Fritz2020}
T.~Fritz.
A synthetic approach to Markov kernels.
\textit{Advances in Mathematics} 370, 2020.

\bibitem{FPR2021}
T.~Fritz, P.~Perrone, S.~Rezagholi.
Three monads on Top and the support as a morphism.
\textit{MSCS} 31(8), 2021.

\bibitem{GKOS2021}
M.~Gaboardi, S.~Katsumata, D.~Orchard, T.~Sato.
Graded Hoare Logic and its Categorical Semantics.
\textit{ESOP}, 2021.

\bibitem{GKLF2020}
N.~Ghani, C.~Kupke, A.~Lambert, F.~Nordvall~Forsberg.
Compositional Game Theory with Mixed Strategies.
\textit{APLAS}, 2020.

\bibitem{KGSU2026}
S.~Kura, M.~Gaboardi, T.~Sekiyama, H.~Unno.
A Category-Theoretic Framework for Dependent Effect Systems.
\textit{ESOP}, 2026.

\bibitem{KW2021}
D.~Kidney, N.~Wu.
Algebras for Weighted Search.
\textit{ICFP}, 2021.

\bibitem{Leinster2021}
T.~Leinster.
\textit{Entropy and Diversity: The Axiomatic Approach}.
Cambridge, 2021.

\bibitem{LCS2025}
J.~Liell-Cock, S.~Staton.
Compositional Imprecise Probability.
\textit{POPL}, 2025.

\bibitem{SK2023}
T.~Sato, S.~Katsumata.
Divergences on Monads for Relational Program Logics.
\textit{MSCS}, 2023.

\end{thebibliography}

\end{document}
